\subsection{位数4の重複固有値1の残存共鳴枝}
\label{subsec:cubic-fourth-root-unit-repeat-resonance}
\begin{theorem}[重複1の最後の位数4二スペクトルも空]
\label{thm:cubic-fourth-root-unit-repeat-resonance-empty}
必要関係重複度 $(r_0,r_1,r_2,r_3)=(2,1,2,1)$ を持つ
半単純スペクトル $(1,1,i,-1)$ と $(1,1,-i,-1)$ のH3包絡は存在しない。
この関係パターンを持つAS4整域自体が存在せず、単冪部分を付けても同様である。
\end{theorem}
\begin{proof}
まず中山が半単純で $\mu=(1,1,i,-1)$ とする。
$U=\langle x,y\rangle$、残る生成元を $b,c$ とする。
右巡回 $\tau$ に対し $w=-(\mu^{-1}\otimes1^{\otimes3})\tau w$ を解く。
解空間は16次元であり、その右巡回軌道代表を次の順にとる：
\[
 xxxy,xxyy,xxcc,xyxy,xyyy,xycc,xbbc,xbcb,
 xcyc,xcbb,xccy,yycc,ybbc,ybcb,ycbb,cccc.
\]
各代表の係数を $a_0,\ldots,a_{15}$ とする。
以下のmiddle blockはこの巡回式から直接得られる。
重み $1$ の語順を $(xx,xy,yx,yy,cc)$、重み $-1$ を
$(xc,yc,bb,cx,cy)$、重み $i$ を $(xb,yb,bx,by)$、重み $-i$ を $(bc,cb)$ とすると
\[
 H_0=\begin{pmatrix}
 0&a_0&-a_0&a_1&a_2\\
 a_0&a_3&-a_1&a_4&a_5\\
 -a_0&-a_1&-a_3&-a_4&-a_{10}\\
 a_1&a_4&-a_4&0&a_{11}\\
 a_2&a_5&-a_{10}&a_{11}&a_{15}
 \end{pmatrix},\quad
 H_2=\begin{pmatrix}
 0&a_8&a_9&-a_2&a_{10}\\
 -a_8&0&a_{14}&-a_5&-a_{11}\\
 -a_9&-a_{14}&0&-a_6&-a_{12}\\
 a_2&a_5&a_6&0&a_8\\
 -a_{10}&a_{11}&a_{12}&-a_8&0
 \end{pmatrix},
\]
\[
 H_{1,3}=\begin{pmatrix}a_6&a_7\\a_{12}&a_{13}\\
 ia_7&ia_9\\ia_{13}&ia_{14}\end{pmatrix}.
\]
Frobenius pairingからrankはそれぞれ $2,2,1$。
$H_{1,3}$ の列空間は一つの関係、行空間は $bc,cb$ 内の一つの関係を与える。
整域性から後者の両係数は非零なので、ある $q\ne0$ に対して第二列は第一列の $q$ 倍。
従って
\[
 a_7=qa_6,\ a_9=q^2a_6,\quad
 a_{13}=qa_{12},\ a_{14}=q^2a_{12}.
\]
非零の $U$-covector $(a_9,a_{14})$ に適合した基底と全体scalarを選び、
$a_9=1,a_{14}=0$ と正規化する。
すると $a_6=q^{-2},a_7=q^{-1},a_{12}=a_{13}=0$。
$H_2$ の五つのPfaffianを零と置くと
\[
 a_8=q^2a_5,\qquad a_{10}=q^4a_5,\qquad a_{11}=0.
\tag{*}
\]
特にその $bb$ 列は常に
\[
 xc+q^{-2}cx\in I_2
\tag{**}
\]
を与える。

$H_0$ の純 $U^2$ 主blockを $A$ とする。
rank $A\le2$ の三次minorには $2a_0^2a_1,2a_1a_4^2$ と
$-a_1(2a_0a_4-a_1^2)$ がある。
$a_1\ne0$ なら前二式が $a_0=a_4=0$ を強制し、第三式に反する。
従って $a_1=0$ である。
$e=(0,1,-1,0)^t$、$v=(a_0,a_3/2,a_3/2,a_4)^t$ とおけば
\[
 A=ev^t+ve^t.
\]
$A=0$ なら、$H_0$ のrank2のため $U^2$ 対 $cc$ の列は非零となる。
しかしその転置側の列が純 $cc$ 関係を与え、整域性に反する。
よって $v\ne0$、$A$ のrankは2。
また $I_0\cap U^2$ は一次元である。
実際、二元二次関係が二本あればそのpencilの行列式は代数閉体上で零点を持ち、
rank1関係が生じる。$I_0$ は二次元なので $cc$ 射影は一次元でなければならない。

$a_4\ne0$ と仮定する。
$(*)$ と $H_0$ の $yy$ 列から純関係 $xy-yx$ を得る。
さらに $U^2$ 対 $cc$ の列 $(a_2,a_5,-q^4a_5,0)^t$ は
$\langle e,v\rangle$ に属するため、$v$ の $yy$ 成分が非零なことから $e$ の倍数であり、
$a_2=0,a_5(1+q^4)=0$。
$a_5=0$ なら $cc$ 列は純 $cc$ 又は零となり、いずれも整域性又は $cc$ 射影の次元に反する。
従って $q^4=-1$。
二本の $I_0$ 関係から $c^2\in\langle U\rangle_2$ であり、$xy=yx$ だから $c^2x=xc^2$。
一方 $(**)$ から $c^2x=q^4xc^2=-xc^2$。整域性と標数0に反する。

従って $a_4=0$ である。
このとき $H_0$ の $yy$ 行は零であり、全二次関係を二側ideal $(x,c)$ で商にすると、
残る関係は $b^2$ 一本だけとなる。
$H_2$ の $bb$ 行は非零なので、その係数も確かに非零である。
従って二次提示そのものから
\[
 B/(x,c)\simeq k\langle y,b\rangle/(b^2).
\]
右辺の次数15は、隣接 $bb$ を持たない語を数えて $1597$ 次元であるが、
AS4型14641の $\dim B_{15}=\binom{18}{3}=816$ より大きい。矛盾する。

$i$ を $-i$ に置き換えた全恒等式は同じ証明を与える。
これはGaloisラベルを $k$ 線型同値と同一視した議論ではない。
最後に一般の $\mu=\mu_s\mu_u$ に対し $H=\mu_u^{1/4}$ をとる。
既証の単冪untwist公式により、右Zhang twist $D=B^H$ は
$\mu_D=\mu_BH^{-4}=\mu_s$ を持つAS4整域である。
$H$ は $\mu_s$ と可換で、二次関係の重み重複度を保つ。
従って上の半単純の場合の矛盾が適用できる。
ここではAS4整域自体の非存在を用いるため、三次 $G$ の正常性をtwistで輸送する仮定は不要である。
\end{proof}

\begin{corollary}[三固有値の九個の必要スペクトル]
\label{cor:cubic-three-eigenvalue-9-list}
三固有値H3の必要リストは位数3の3個、位数4の6個、合計9個、Galoisラベル5軌道へ縮小する。
残存候補の存在・許容性及び全H3の分類は未完成である。
\end{corollary}
